\documentclass{plantillaPPT}
\titulo{9}{Integrales sobre regiones no rectangulares}
\begin{document}
\primeraDiapo
\begin{frame}{Problema 1}
\framesubtitle{Práctica 9}
1. Calcular el volumen del sólido acotado por la superficie $z = \sin(y)$, los planos $x=0$, $x=1$, $y=0$, $y=\pi/2$ y el plano $xy$.
\end{frame}
\begin{frame}{Problema 2}
\framesubtitle{Práctica 9}
2. Si $f(x,y)=e^{\sin(x+y)}$ y $R=[-\pi,\pi]\times[-\pi,\pi]$. Mostrar que
\begin{equation*}
    \frac{1}{e}\leq\frac{1}{4\pi^2}\iint_R e^{\sin(x+y)}\,dA \leq e
\end{equation*}
\end{frame}
\begin{frame}{Problema 3}
\framesubtitle{Práctica 9}
3. $(a)$ Hallar la integral de 
\begin{equation*}
    \int_0^4\int_{x/2}^2e^{y^2}\,dydx
\end{equation*}
$(b)$ Dado que la integral doble $\iint_R f(x,y)\,dA$ de una función continua positiva $f$ es igual a la integral iterada
\begin{equation*}
    \int_0^1 \left[\int_y^{\sqrt{2-y^2}}f(x,y)\,dx\right]\,dy
\end{equation*}
Esbozar la región R e intercambiar el orden de integración.
\end{frame}
\begin{frame}{Problema 4}
\framesubtitle{Práctica 9}
4. Hallar el volumen de la región determinada $x^2+y^2+z^2\leq10$, $z\geq2$.
Usar el método del disco (Cálculo 2), y decir como está relacionado con el principio de Cavalieri.
\end{frame}
\end{document}